\section*{Definiciones y Medidas de Seguridad}

% Corrección de términos y ortografía
Debido a que en las actividades normales de la radiografía industrial es frecuente trabajar en áreas imposibles de evacuar, es necesario tomar una serie de medidas tendientes a proteger al público o a trabajadores no ocupacionalmente expuestos.

% Definición de zonas y medidas específicas
Dentro de estas medidas se incluyen la delimitación de áreas y su señalización adecuada. Principalmente se busca definir dos tipos de zonas de acuerdo a lo siguiente:

\subsection*{Zona Restringida (Zona de Radiación)}
Es aquella a la que por motivos de seguridad radiológica se prohíbe el acceso o permanencia a todo el personal excepto al técnico encargado de la fuente o a su(s) auxiliar(es) y siempre que dicha fuente no se encuentre fuera de su contenedor.

% Definición de niveles de radiación permitidos
Toda zona con niveles de exposición iguales o superiores a $5 \text{ mR/h}$ será considerada como zona restringida. La zona restringida se mantendrá acordonada y debidamente señalizada con letreros indicando la presencia de radiación y con distancias entre dos consecutivos no mayores a $10$ metros.

\subsection*{Zona Controlada}
En las áreas donde el nivel de radiactividad es controlado, se vigila el flujo de personal. Se considera como zona comprendida en la restringida pero con un nivel de exposición mayor a $2 \text{ mR/h}$. La zona controlada no requiere acordonamiento pero se mantendrá vigilada por el personal responsable a fin de evitar que alguna persona pudiera permanecer innecesariamente en ese lugar.

% Establecimiento de límites máximos de dosis de radiación
Los límites de dosis administrados de radiación serán los indicativos limites operativos de Indesa en nuestro manual de Seguridad Radiológica son los siguientes:
\begin{itemize}
  \item El P.O.E. no deberá recibir una dosis mayor a $4 \text{ mSv} (400 \text{ mrem})$ por mes.
  \item El P.O.E. no deberá rebasar $5 \text{ mSv} (5000 \text{ mrem})$ en un año.
\end{itemize}

% Cumplimiento con los niveles de referencia y de registro
Y se deberá cumplir con los niveles de referencia y de registro indicados en nuestro manual de Seguridad Radiológica parte 3.2.

% Comentario sobre la efectividad de las restricciones
Las restricciones sobre las dosis efectivas son suficientes para asegurar que se eviten la afectación detrimental en los tejidos y órganos del cuerpo, excepto el cristalino del ojo, el cual hace una contribución despreciable a la dosis efectiva, y la piel que puede estar sujeta a exposiciones localizadas. Se mencionan límites especiales para estas tecjidos. Los límites anuales son $150 \text{ mSv}$ para el cristalino y $500 \text{ mSv}$ para la piel, promediados sobre cualquier centro cuadrado independientemente del area expuesta.
